\documentclass[10pt,a4paper]{amsart}
\usepackage[margin=19mm]{geometry}
\usepackage{amsmath,amssymb,mathtools}
\usepackage[scr=rsfs,cal=euler]{mathalfa}
\usepackage{xfrac}
\usepackage[unicode,hidelinks,bookmarks=false]{hyperref}
\newcommand{\ie}{i.e.\ }
\newcommand{\cf}{cf.\ }
\newcommand{\resp}{respectively\ }
\newcommand{\Subsection}{\subsection}
\providecommand{\nobreakdash}{\nobreak\hbox{-}}
\setlength{\emergencystretch}{2em}
\multlinegap0pt
\usepackage{amssymb}
\usepackage[shortlabels]{enumitem}
\newtheorem{lemma}{Lemma}
\newcommand{\A}{\mathbf{A}}
\newcommand*\Div{\operatorname{div}}
\newcommand{\R}{\mathbb{R}}
\title{On the nondivergent variant of the existence result of Boccardo and Orsina}
\author{Agnieszka Ka{\l}amajska, Dalimil Pe\v{s}a, Artur Rutkowski}
\date{Source: arXiv:2605.03704v1 (CC BY 4.0)}
\begin{document}
\maketitle

\noindent\emph{Source and licence.} On the nondivergent variant of the existence result of Boccardo and Orsina. Version arXiv:2605.03704v1 (CC BY 4.0), distributed by arXiv under the Creative Commons Attribution 4.0 International licence, http://creativecommons.org/licenses/by/4.0/, as recorded in the arXiv licence metadata for this exact version and shown on the abstract page https://arxiv.org/abs/2605.03704v1. Full licence notice: \texttt{licenses/arxiv-2605.03704-CC-BY-4.0.txt}.\par\smallskip

\noindent\textit{Editorial scope.} Counted unit: the article's lemma lem:specialxi (source0.tex lines 967--974). The article's complete original proof of it is delivered with it (source0.tex lines 975--1053, one \texttt{proof} environment of 79 lines). No mathematics is written, repaired or re-derived: the statement and the proof are the article's own lines, in the article's own notation, and no retained line is substituted or re-flowed.  Retained uncounted as the source-local context the proof needs: lines 153--155; lines 823--825; lines 897--899.  Nothing was omitted from any retained span, so the package carries no omission record.  No retained line contains a citation, so the wrapper carries no bibliography and no bibliography entry.  No retained line contains a figure environment, an image-inclusion command or drawing code, so no image is delivered and the package declares no asset dependency.  Editorial formatting only, in addition: an AMS \texttt{amsart} preamble that restates the article's own declarations and macros used by the retained lines, namely the environments \texttt{lemma} on separate counters, together with the standard packages those lines need.  The retained environments print in this document as Lemma 1 in order of appearance, whereas the article prints the same items with the numbering of its own full document.  The complete native source of the article as distributed in the arXiv e-print for this exact version is bundled unchanged as \texttt{sources/199812/source0.tex}.
\begin{align} \tag{$\A$} \label{ellipticity}
 c_{\A}\|\xi\|^2 \le \A(x)\xi \cdot \xi \le C_{\A} \|\xi\|^2,\quad x\in \overline{\Omega},\ \xi\in \R^d \quad \text{ for some } 0< c_{\A} \leq C_{\A} <\infty.
\end{align}

\begin{align} \label{eq:P*non-div-P}
    P^*v = Pv + 2\Div \A \cdot \nabla v + (\Div(\Div \A))v,
\end{align}

\begin{align}\label{eq:firstextension}
    \mathcal{A} = \{\phi\in C^1_c(\Omega)\cap W^{2,1}(\Omega): P^*\phi\in L^\infty(\Omega)\}.
\end{align}

\begin{lemma}[approximation of $\xi$]\label{lem:specialxi}
    Assume that $\Omega \in C^{2}$ and that $\A\in C^{2}(\overline{\Omega})$ satisfies \eqref{ellipticity}. Then there exists an increasing sequence of positive functions $\xi_n\nearrow \xi$ such that $\xi_n \in \mathcal{A}$ (see \eqref{eq:firstextension}) 
    and for almost every $x\in \Omega$
    $$
        -P^*\xi_n (x) \le D_{\A},
    $$
    where  $D_{\A}=  1+ \| (\Div(\Div\A ))_+\|_{\infty}\|\xi\|_\infty$.
\end{lemma}

\begin{proof}
    \smallskip
    \noindent
    Let us start by the construction of $\xi_n$.
    
     Consider a convex function $\Phi\in C^\infty (\mathbb{R})$ such that $\Phi$ is nondecreasing, $\Phi(x)= 0$ for $x<\delta$ for some $\delta>0$, $\Phi$ is Lipschitz with constant $1$ (in particular  $\Phi(t)\le t$ and $\Phi' \leq 1$),  
     %$\Phi\equiv 0$ in some neighborhood of zero and $\Phi(t)\le t$ 
     and $\frac{\Phi(t)}{t}\nearrow 1$ as $t\to \infty$.

     As an example function $\Phi$ we can consider $\Phi (t):= \varphi_\epsilon * {\rm max}\{ 0, t-1\}$, where $\varphi_\epsilon$ is the standard mollifier.

    Define 
    $$
        \xi_n(x):= \frac{1}{n}\Phi (n\xi (x)) .
    $$
    By the regularity and decay of $\xi$ and the definition of $\Phi$ we have $\xi_n\in C^1_c(\Omega)$. 
    Furthermore, 
    $$
        \xi_n (x) = \frac{\Phi (n\xi (x))}{n\xi (x)}\xi (x) \le \xi(x).
    $$
    and $\Phi(n\xi)/n\xi$ converge to 1 as $n\to \infty$, and increase since $\Phi$ is convex. This implies that 
    $
        \xi_n(x)\nearrow \xi (x)
    $
    for every $x\in \Omega$.
    
    
    We will show that
     the following estimate holds almost everywhere:
    \begin{align*}
    %\Div( \A\nabla\xi_n)&\le \Phi^{'}(n\xi);\nonumber\\
    - P^*\xi_n &\le \Phi^{'}(n\xi) + \Div(\Div\A)(\Phi^{'}(n\xi)\cdot\xi-\frac{1}{n}\Phi(n\xi)).
    \end{align*}
    
    Indeed, $\xi_n\in W^{2,1}_c(\Omega)$ and
    \begin{align}
        \nabla \xi_n   &=\Phi^{'}(n\xi )\nabla\xi  ,\nonumber\\
        \nabla^{(2)}\xi_n  &= \Phi^{'}(n\xi )\nabla^{(2)}\xi  + n\Phi^{''}(n\xi )\nabla\xi  \otimes \nabla\xi  ,\nonumber\\
        \Div(\A\nabla\xi_n) &= (\Div\A)\cdot \nabla\xi_n + \A\cdot \nabla^{(2)}\xi_n\label{1oszac}\\
        &= \Phi^{'}(n\xi )(\Div\A)\cdot \nabla\xi   + \Phi^{'}(n\xi ) \A\cdot\nabla^{(2)}\xi  \nonumber\\ 
        & \quad + n\Phi^{''}(n\xi ) \A \nabla\xi  \cdot \nabla\xi  \nonumber\\
        &\ge  \Phi^{'}(n\xi )\Div(\A \nabla\xi ) ,\nonumber
    \end{align}
    where we have used the convexity of $\Phi$ and \eqref{ellipticity}, which imply $\Phi^{''}(n\xi )(\A \nabla\xi \cdot\nabla \xi)\ge 0$.
    We also have
    \begin{align}\label{eq:divdivterm}
    \Div(\Div\A\xi_n) &= \Div(\Div\A)\xi_n + (\Div\A)\cdot\nabla \xi_n\\
    &= \Div(\Div\A)\frac{1}{n}\Phi(n\xi ) + 
     \Phi^{'}(n\xi )(\Div\A)\cdot \nabla\xi  \nonumber. 
    \end{align}
    Consequently,
    \begin{align*}
    -P^*\xi_n &= -\Div(\A\nabla\xi_n) - \Div(\Div\A\xi_n)\\
    &\! \!\stackrel{\eqref{1oszac} }{\le}  - \Phi^{'}(n\xi )[\Div(\A\nabla\xi )]  -  \Phi^{'}(n\xi ) \Div\A\cdot \nabla\xi   - \Div(\Div\A)\frac{1}{n}\Phi(n\xi )\\
    &=  -\Phi^{'}(n\xi )[ \Div(\A \nabla\xi ) + \Div(\Div\A \xi )]  \\ &\quad -  \Phi^{'}(n\xi )[\Div\A\cdot \nabla\xi  
    -\Div(\Div\A \xi )]\\
    &\quad - \Div(\Div\A)\frac{1}{n}\Phi(n\xi )\\
    &= \Phi^{'}(n\xi )[-P^*\xi] + \Phi^{'}(n\xi )
    [\Div(\Div\A ) \xi] -  \Div(\Div\A)\frac{1}{n}\Phi(n\xi )\\
    &=  \Phi^{'}(n\xi ) +   \Div(\Div\A )\left( \Phi^{'}(n\xi )\, \xi - \frac{1}{n}\Phi(n\xi )\right) \\
    &\le 1+ \| (\Div(\Div\A ))_+\|_{\infty}\|\xi\|_\infty.
    \end{align*}
    The last estimate holds because
    \begin{equation*}
        \frac{1}{n\xi}\Phi(n\xi ) = \frac{\Phi(n\xi ) - \Phi(0)}{n\xi} = \Phi'(\zeta)
    \end{equation*}
    for some $\zeta \in (0, n\xi)$ by the Lagrange Mean Value Theorem, whence the convexity of $\Phi$ implies 
    \begin{equation*}
        0 \leq \Phi^{'}(n\xi ) - \frac{1}{n \xi}\Phi(n\xi ) \leq \Phi^{'}(n\xi ) \leq 1.
    \end{equation*}
    It remains to show that $\xi_n\in \mathcal{A}$. By recalling \eqref{eq:P*non-div-P} and applying \eqref{1oszac} and \eqref{eq:divdivterm} we have
    \begin{equation*}
    \begin{split}
         P^\ast \xi_n &= 
         \Phi'(n\xi) P^\ast \xi + n\Phi^{''}(n\xi ) \A \nabla\xi  \cdot \nabla\xi + \Div (\Div \A) (\xi_n - \Phi'(n\xi) \xi)
    \end{split}
    \end{equation*}
    As $P^\ast \xi = 1 \in L^{\infty}(\Omega)$, $0\leq \Phi' \leq 1$, $\Phi \in C^{\infty}(\mathbf{R})$, $\xi \in C^1(\overline{\Omega})$, $\A \in C^2(\overline{\Omega})$ satisfies \eqref{ellipticity}, and $\xi_n \in C^1_c(\overline{\Omega})$, we obtain $P^\ast \xi_n \in L^{\infty}(\Omega)$, hence $\xi_n \in \mathcal{A}$, which ends the proof.
\end{proof}

\end{document}
